% Vietnamese translation for OI Public Library
% Source-File: IOI中国国家候选队论文/2006/李天翼/李天翼.ppt
\documentclass[11pt]{article}
\usepackage{oipl-vn}

\title{Xét từ các trường hợp đặc biệt\\{\large Ghi chú theo slide}}
\author{Li Tianyi}
\date{2006}

\begin{document}
\maketitle

\origtitle{Reasoning from special cases}
\sourcepath{IOI China candidate team papers / 2006 / Li Tianyi / Li Tianyi.ppt}

\section{Phạm vi}

Bộ slide đi cùng bài viết về cách phân tích trường hợp đặc biệt. Ghi chú này
dựa trên bản DOC/PDF đã dịch, trong đó các ví dụ chính là \emph{Bra},
\emph{Sko} và \emph{Polygon}.

\section{Hai loại trường hợp}

Trường hợp đặc biệt thường thuộc hai nhóm. Nhóm thứ nhất là trường hợp đơn
giản, nơi một vài tham số nhỏ hoặc quan hệ bị thu hẹp làm bản chất bài toán
lộ rõ. Nhóm thứ hai là trường hợp cực đoan, nơi nghiệm tối ưu bị ép vào một
dạng biên. Cả hai nhóm đều giúp ta tìm bất biến, trạng thái hoặc phép biến
đổi cho bài tổng quát.

\section{Từ ví dụ nhỏ tới mô hình}

Khi một bài có nhiều điều kiện rối, thử giải tay dữ liệu nhỏ thường giúp nhận
ra điều gì thật sự quan trọng. Với mạch trong \emph{Bra}, trạng thái cố định
của cổng không nên được suy luận từng cổng độc lập; ta cần nhìn quan hệ giữa
các trạng thái có thể cùng tồn tại trong toàn mạch.

\section{Từ biên tới thuật toán}

Trong các bài hình học hoặc tổ hợp, trường hợp cực đoan thường cho biết nghiệm
tối ưu chạm vào một biên nào đó: một đỉnh, một cạnh, một thứ tự đã sắp xếp,
hoặc một cấu hình suy biến. Sau khi xác định được biên này, thuật toán có thể
duyệt ít khả năng hơn nhiều so với không gian ban đầu.

\section{Ghi nhớ}

Trường hợp đặc biệt không phải là phần phụ của bài toán. Nó là công cụ để kiểm
tra giả thuyết và tìm cấu trúc. Một lời giải tốt thường bắt đầu từ quan sát ở
trường hợp nhỏ, rồi chứng minh quan sát ấy vẫn đúng trong bài tổng quát.

\end{document}
